\newpage
\pagestyle{plain}
\section{\centering ВВЕДЕНИЕ}
Численное решение задачи Коши для обыкновенных дифференциальных уравнений является одной из фундаментальных математических задач. Однако во многих случаях аналитическое решение таких задач невозможно \cite{Verzhbitsky}, поэтому возникла необходимость разработки и исследования различных математических моделей и численных методов. 

Наибольшее распространение среди методов получили методы Рунге-Кутты вследствие их высокой точности и универсальности. В реальных задачах часто применяются методы высоких порядков точности, построение которых связано с выполнением системы нелинейных условий на коэффициенты метода. Несмотря на то, что методы Рунге-Кутты изучаются давно, нахождение коэффициентов методов в явном виде затруднено вследствие существенного возрастания размера и сложности системы \cite{Demidovich}. 

Традиционные подходы к выбору коэффициентов, которые основаны на анализе локальных характеристик метода, являются недостаточными для ряда задач, в том числе для задач с осциллирующими решениями, вследствие того, что они не учитывают поведение численного решения на всём интервале интегрирования. Более того, в осциллирующих задачах стандартные методы демонстрируют численную ошибку, несообразную порядку точности \cite{article2023}. Поэтому возникает необходимость использования новых подходов для подбора коэффициентов.

Один из таких подходов заключается в рассмотрении семейства параметрических методов Рунге-Кутты, т.е. методов, содержащих свободные коэффициенты. Его преимущество заключается в возможности выбора из возможных схем наиболее подходящей для конкретной задачи, не фиксируя метод заранее. Именно оптимальный подбор параметров обеспечивает эффективность метода. Поэтому применение методов оптимизации для подбора численных схем является перспективным способом повышения точности численных методов за счёт адаптации метода к решаемой задаче \cite{article2023}. В частности, в ряде задач, где исследуются сложные многомерные пространства параметров, используется эволюционный алгоритм, преимущество которого заключается в отсутствии необходимости вычисления производных целевой функции \cite{article2025}.

Актуальность данной работы обусловлена необходимостью повышения точности численного решения задач Коши для обыкновенных дифференциальных уравнений с осциллирующими решениями, а также разработки методов, способных учитывать глобальные свойства решения на всём интервале интегрирования. Использование параметрических семейств методов Рунге-Кутты в сочетании с современными методами оптимизации позволяет существенно улучшить качество численного решения и расширить область применимости классических численных схем.

Целью данной работы является разработка и исследование параметрических методов Рунге–Кутты, ориентированных на эффективное решение задачи Коши с осциллирующими решениями, а также изучение подходов к оптимизации их коэффициентов.

Для достижения поставленной цели необходимо решить следующие задачи:
\begin{itemize}
	\item провести анализ существующих методов Рунге–Кутты и условий их построения;
	\item исследовать особенности численного решения задач с осциллирующими решениями;
	\item построить семейства параметрических методов Рунге–Кутты;
	\item разработать подходы к оптимизации коэффициентов методов, в том числе с использованием эволюционных алгоритмов;
	\item провести аналитическое исследование свойств полученных методов;
	\item реализовать разработанные методы и выполнить их численное тестирование;
	\item провести сравнительный анализ эффективности предложенных методов.
\end{itemize}

Таким образом, в работе рассматривается задача построения и исследования параметрических методов Рунге–Кутты, а также разработки эффективных алгоритмов оптимизации их коэффициентов для повышения точности численного решения задачи Коши.